\hypertarget{appendix-83c3f28f7bda}{}
\section*{Positional notation}
\section{Positional notation}\label{positional-notation}

\textbf{Curriculum depth:} Introduced · \textbf{First seen:} C001 · \textbf{Latest development:} C001

\subsection{Definition and mechanism}\label{definition-and-mechanism}

Digits receive powers-of-base weights; binary uses base two and \(V=\sum b_i2^i\).

A base is the number of available digits. Position i contributes digit × baseⁱ; the rightmost position is i=0. For \texttt{1011₂}, the weights 8,4,2,1 give 8+2+1=11. The number of states is 2ᴺ, while the maximum unsigned value is 2ᴺ−1.

\subsection{Prerequisites}\label{prerequisites}

\begin{itemize}
\tightlist
\item
  \hyperref[appendix-ddbb69998a2e]{Bit} --- A bit has two possible logical states.
\end{itemize}

\subsection{Used by}\label{used-by}

\begin{itemize}
\item
  \hyperref[appendix-2a2f1b74bdd3]{Terminating fractions and dyadic rationals} --- Finite fractional digits correspond to an integer over a power of the base.
\item
  \hyperref[appendix-a24fd710cf67]{Significand and hidden bit} --- uses positional notation to establish this mechanism's representation or policy.
\item
  \hyperref[appendix-dc552784a2c0]{Binary conversion} --- uses positional notation as part of its representation or reasoning.
\item
  \hyperref[appendix-6b8075763724]{Hexadecimal} --- uses positional notation as part of its representation or reasoning.
\item
  \hyperref[appendix-17fc1bc8b305]{Octal} --- uses positional notation as part of its representation or reasoning.
\item
  \hyperref[appendix-280445746493]{Bit significance} --- uses positional notation as part of its representation or reasoning.
\item
  \hyperref[appendix-c99977d88b33]{Unsigned modular arithmetic} --- uses positional notation as part of its representation or reasoning.
\item
  \hyperref[appendix-ce20506c8532]{Signed magnitude} --- uses positional notation as part of its representation or reasoning.
\item
  \hyperref[appendix-d52d5f78ffd1]{Carry and borrow} --- uses positional notation as part of its representation or reasoning.
\item
  \hyperref[appendix-92e040df969a]{Big endian} --- depends on this concept; see its prerequisite explanation.
\end{itemize}

\subsection{Recall and next depth}\label{recall-and-next-depth}

Explain positional notation without the label, work the example above, then state the assumption whose removal would change the result.

\textbf{Pending:} signed and fractional encodings.

\subsection{Source and navigation}\label{source-and-navigation}

\href{Continuous\%20Notes/Computing\%20Foundations\%20-\%20From\%20Bits\%20to\%20C++\%20and\%20Systems\%20APIs\#From\%20a\%20signal\%20to\%20an\%20interpretation}{C001 teaching chapter} · \hyperref[appendix-fd090d982a15]{Part definitions and mechanisms} · \href{}{Dependency Map (02-Foundations-II, appendix-4ad47bdd269a)}
